\documentclass[12pt]{../../oss-apphys-exam}
\usepackage{enumitem}

\usetikzlibrary{patterns}

\tikzset{
  >=latex
}

\runningheader{
  \fontfamily{phv}\small
  \textcolor{gray}{AP PHYSICS C: MECHANICS $\blacksquare$}
  \textbf{FREE-RESPONSE QUESTIONS}}{}{}


\begin{document}

\begin{center}
  {
    \fontfamily{phv}\Large
    \textbf{PHYSICS C: MECHANICS\\SECTION II\\TIME -- 60 MINUTES}
  }
\end{center}

\vspace{2in}
\textbf{Directions:}

Section II has 3 questions and lasts 60 minutes.

You may use the available paper for scratch work and planning, but you must
write your answers in the free-response booklet. Label parts (e.g., \textbf{A},
\textbf{B}, \textbf{C}) and sub-parts (e.g., i, ii, iii) as needed. Use a
pencil or a pen with black or dark blue ink to write your responses.

A calculator is allowed in this section, as well as a ruler and straightedge.
You may use a handheld four-function, scientific, or graphing calculator, or
the calculator available in this application. Reference information, including
lists of equations, can also be accessed in this application and is available
throughout the exam.

All final numerical answers should include appropriate units when applicable.
Credit for your work depends on demonstrating that you know which physical
principles to apply in a particular situation. Credit will be awarded only for
work that is clearly designated as the solution to a specific part of a
question. Credit also depends on the quality of your solutions and explanations.
Therefore, you should show your work for each part in the space provided for
that part. If you need more space, be sure to clearly indicate where you
continue your work. When constructing a graph or diagram, use only one color of
ink or pencil.

You may pace yourself as you answer the questions in this section, or you may
use these optional timing recommendations:

It is suggested that you spend about 20 minutes on each question.

You can go back and forth between questions in this section until time expires.
The clock will turn red when 5 minutes remain---\textbf{the proctor will not
  give you any time updates or warnings.}
\newpage


% THIS QUESTION IS BASED ON THE 2024 AP PHYSICS C MECHANICS EXAM SET 1
% QUESTION 1. FOR NOW, IT IS USED WITHOUT MODIFICATIONS

\begin{center}
  \pic{.7}{two-blocks-slide}
        
  \underline{Note:} Figure not drawn to scale.
        
  Figure 1
\end{center}

\begin{questions} 
  \question Block A and Block B of masses $m$ and $3m$, respectively, are
  arranged in a setup consisting of an ideal spring with spring constant $k$
  and a horizontal surface. Friction between the surface and the blocks is
  negligible except in a region of length $D$, where the coefficient of kinetic
  friction between Block A and the surface is $\mu$. Block B is attached to a
  string of length $\ell$ and negligible mass, as shown in Figure 1. Block A
  is held against the spring, compressing the spring a distance $x_c$.

  \vspace{.2in}At time $t=0$, Block A is located at position $x = x_0$ and is
  released from rest. After the block is released, the following occurs.

  \begin{itemize}[itemsep=6pt]  
  \item At time $t=t_1$, Block A is at $x=x_1$ after traveling a distance
    $x_c$. Block A moves with speed $v$, and the spring is at its equilibrium
    position.

  \item At time $t=t_2$, the left side of Block A is at $x=x_2$ after passing
    through a distance $D$across the region with nonnegligible friction.

  \item At time $t=t_3$, Block A is at $x=x_3$ and Block A collides with and
    sticks to Block B.
  \end{itemize}

  For parts \textbf{A} and \textbf{B}, express your answer in terms
  of $m$, $k$, $D$, $\mu$, $x_c$, and physical constants, as appropriate.
  \begin{parts}
    \part\textbf{Derive} an expression for the speed $v$ of Block A at
    time $t_1$.
    
    \part\textbf{Derive} an expression for the speed $v_{A,B}$ of the two-block
    system immediately after the collision at time $t_3$.
    
    \part On the following axes, sketch a graph of the kinetic energy $K$ of
    Block A as a function of time $t$ from time $t=0$ to time $t_3$.
    \label{sketch}
    \begin{center}
      \begin{tikzpicture}[scale=.9]
        \draw[axes] (0,0)--(10,0) node[right]{$t$};
        \draw[axes] (0,0)--(0,5.5) node[above]{$K$};
        \foreach \x in {1,2,3}{
          \draw[dashed] (\x*2.5,0)--+(0,5) node[pos=0,below]{$t_\x$};
        }
      \end{tikzpicture}
    \end{center}
    
    \part Use principles of work and energy to \textbf{justify} the graph drawn
    in part \textbf{C} for the time interval $t=0$ to $t=t_1$. Explicitly
    reference features of the shape of the graph you drew in part \textbf{C}.
    \vspace{.5in}
    
    \uplevel{
      \begin{center}
        \pic{.3}{swing}
        
        \underline{Note:} Figure not drawn to scale.
        
        Figure 2
      \end{center}
      After the collision, the two-block system instantaneously comes to rest
      at time $t_4$ which occurs when the string makes a small angle
      $\theta_\text{max}$ with the vertical, as shown in Figure 2. For times
      $t>t_4$, the system oscillates with frequency $f_\ell$. The support holding
      the string is raised, and the procedure is then repeated using a new
      string of length $2\ell$
    }

    \part\textbf{Indicate} how the new frequency of oscillation $f_{2\ell}$ of
    the system on the new string of length $2\ell$ will compare to the
    frequency of oscillation $f_\ell$ from the original procedure.

    \vspace{.2in}
    \underline{\hspace{.5in}} $f_\ell > f_{2\ell}$\hspace{.5in}
    \underline{\hspace{.5in}} $f_\ell < f_{2\ell}$\hspace{.5in}
    \underline{\hspace{.5in}} $f_\ell = f_{2\ell}$
    
    \vspace{.2in}Briefly \textbf{justify} your answer.
  \end{parts}
  \newpage


  % THIS QUESTION IS BASED ON THE 2000 AP PHYSICS C MECHANICS EXAM QUESTION 1.
  % FOR NOW, IT IS USED WITHOUT MODIFICATIONS
  \question A rubber ball of mass $m$ is dropped from a cliff. As the ball
  falls, it is subject to air drag (a resistive force caused by the air). The
  drag force on the ball has magnitude $bv^2$, where $b$ is a constant drag
  coefficient and $v$ is the instantaneous speed of the ball. The drag
  coefficient $b$ is directly proportional to the cross-sectional area of the
  ball and the density of the air and does not depend on the mass of the ball.
  As the ball falls, its speed approaches a constant value called the terminal
  speed.
  \begin{parts}
    \part On the figure below, \textbf{draw} and \textbf{label} all the forces
    on the ball at some instant before it reaches terminal speed. The forces
    should be drawn as arrows originating at, and pointing away from, the
    figure.
    \begin{center}
      \vspace{.5in}{\tikz\draw[very thick,fill=lightgray] circle (.13);}
      \vspace{.5in}
    \end{center}
    
    \part\textbf{Indicate} whether the magnitude of the acceleration of the
    ball of mass $m$ increases, decreases, or remains the same as the ball
    approaches terminal speed. \textbf{Justify} your answer.
    
    \vspace{.15in}
    \underline{\hspace{.3in}} Increases\hspace{.4in}
    \underline{\hspace{.3in}} Decreases\hspace{.4in}
    \underline{\hspace{.3in}} Remains the same
    
    \part\textbf{Write}, but do NOT solve, a differential equation for the
    instantaneous speed $v$ of the ball in terms of time $t$, the given
    quantities, and fundamental constants.
    
    \part\textbf{Determine} the terminal speed $v_t$  in terms of the given
    quantities and fundamental constants.
    
    \part\textbf{Determine} the energy dissipated by the drag force during the
    fall if the ball is released at height $h$ and reaches its terminal speed
    before hitting the ground, in terms of the given quantities and fundamental
    constants.
  \end{parts}
  \newpage


  \question A nonlinear spring is compressed various distances $x$, and the
  force $F$ required to compress it is measured for each distance. The data are
  shown in the table below.
  \begin{center}
    \begin{tabular}{|c|c|p{1in}|}
      \hline
      $x$ [m] & $F$ [N] & \\ \hline
      0.05 & 4 & \\ \hline
      0.10 & 17 & \\ \hline
      0.15 & 38 & \\ \hline
      0.20 & 68 & \\ \hline
      0.25 & 106 & \\ \hline
    \end{tabular}
  \end{center}
  Assume that the magnitude of the force applied by the spring is of the form
  $F(x)=Ax^2$.
  \begin{parts}
    \part\textbf{Explain} which quantities should be graphed in order to yield
    a straight line whose slope could be used to calculate a numerical value
    for $A$.
    
    \part\textbf{Calculate} values for any of the quantities identified in part
    \textbf{A} that are not given in the data, and record these values in the
    table above. Label the top of the column, including units.
    
    \part On the axes below, \textbf{plot} the quantities you indicated in part
    \textbf{A}. \textbf{Label} the axes with the variables and appropriate
    numbers to indicate the scale.
    \uplevel{
      \centering
      \begin{tikzpicture}[scale=.7]
        \draw[axes] (0,0)--(14.5,0);
        \draw[axes] (0,0)--(0,11);
        \draw[dashed,gray] grid (13.8,10.8);
      \end{tikzpicture}
    }
    
    \part Using your graph, calculate $A$.
    
    \uplevel{
      The spring is then placed horizontally on the floor. One end of the
      spring is fixed to a wall. A cart of mass 0.50 kg moves on the floor with
      negligible friction and collides head-on with the free end of the spring,
      compressing it a maximum distance of \SI{.10}\metre.
    }
    
    \part\textbf{Calculate} the work done by the cart in compressing the spring
    \SI{.10}{\metre} from its equilibrium length.
    
    \part\textbf{Calculate} the speed of the cart just before it strikes the
    spring.
  \end{parts}    
\end{questions}
\end{document}
